% VCITE Citation Example — LaTeX
% Demonstrates the \vcite macro syntax for embedding passage-level
% verifiable citations. Requires a vcite.sty package (not provided
% here; this file shows the macro interface only).

\documentclass[12pt]{article}
\usepackage[utf8]{inputenc}
\usepackage[T1]{fontenc}
\usepackage{hyperref}
\usepackage{xcolor}

% ── VCITE Macro Definitions (inline; normally from vcite.sty) ──

% \vcite[relation]{key} — inline marker, no visible text change
%   Writes a \label and stores metadata for the VCITE manifest.
%   The key matches a .bib entry; the passage is taken from the
%   surrounding \begin{vcitepassage}...\end{vcitepassage} environment.

\newcommand{\vcite}[2][supports]{%
  \textsuperscript{\textcolor{gray}{\tiny\textsf{V}}}%
  % In production, vcite.sty writes to \jobname.vcite aux file:
  %   \vcite@entry{#2}{relation=#1}
}

% \begin{vcitepassage}{id}{hash}
%   Wraps the cited passage. The hash is the VCITE SHA-256 fingerprint.
%   Content between \begin and \end is the text_exact.
% \end{vcitepassage}

\newenvironment{vcitepassage}[2]{%
  % #1 = vcite-id, #2 = sha256 hash
  % In production, vcite.sty highlights the passage and records
  % it in the .vcite manifest for post-processing.
  \ignorespaces
}{%
  \ignorespacesafterend
}

% ── Document ──────────────────────────────────────────────

\begin{document}

\title{The Link Rot Problem in Scholarly Communication}
\author{VCITE Example}
\date{}
\maketitle

\section{Introduction}

Digital scholarship depends on stable references, but the web's
mutability undermines this assumption. As documented by Zittrain
et~al., %
\begin{vcitepassage}{vcite-d2ba5887}{sha256:d2ba5887d4456199ce37a7860ce011485519925dc75c3564f50d1ffa1e1dc933}%
over 70\% of the URLs provided in articles published in the Harvard
Law Review did not lead to originally cited information%
\end{vcitepassage}%
\vcite[supports]{zittrain2014perma}.
This rate of decay accelerates over time, raising questions about the
long-term integrity of the scholarly record.

The Columbia Journalism Review found that %
\begin{vcitepassage}{vcite-26c71181}{sha256:26c71181069e54c4d776f2a03de838c81aab3eff1a64a71ecde920f629367707}%
eight major AI search tools provided incorrect or inaccurate answers
to more than 60\% of 1,600 test queries%
\end{vcitepassage}%
\vcite[quantifies]{lim2025aisearch},
with error rates ranging across tools and query types.

\section{VCITE Macro Reference}

The \verb|\vcite| macro and \texttt{vcitepassage} environment work
together:

\begin{enumerate}
  \item \verb|\begin{vcitepassage}{id}{hash}| wraps the verbatim
        cited text. The \texttt{id} is the VCITE object identifier
        and the \texttt{hash} is the SHA-256 passage fingerprint.
  \item \verb|\vcite[relation]{bibkey}| places a superscript marker
        and links the passage to a bibliography entry. The optional
        \texttt{relation} argument defaults to \texttt{supports}.
  \item At \verb|\end{document}|, \texttt{vcite.sty} writes a
        \verb|\jobname.vcite| manifest (JSON) listing all VCITE
        objects with their passage text, hashes, and relations.
\end{enumerate}

Available relation values: \texttt{supports}, \texttt{contradicts},
\texttt{defines}, \texttt{quantifies}, \texttt{contextualizes},
\texttt{method}, \texttt{cautions}, or any \texttt{x-*} extension.

\end{document}
